\cotefile{157-4}

\scanpage{2}

\(\underline{\overline{\mathrm{XIII}}}\) Construction de \unsure{catègories|catégories} de \unsure{modèles|modeles} et de dérivateurs

(1)

(\unsure{Variances|Variance} $\unsure{u^{\ast}_{!}|U^{\ast}_{!}|W^{\ast}_{!}|u^{?}_{!}},\unsure{u^{\ast}|u'},\unsure{u_{\ast}|u_{\times}}$ p.~44--50)

\unsure{def.|déf.} parties $\unsure{\pi|\Pi}$-\unsure{accessibles} page \unsure{28|88}

% archive sheet number: 1

\scanpage{3}

\mspage{1}

Catégories de modèles.

1. \deleted{\unsure{2|3}}\uline{Couples de Quillen.} $(\Phi,\Psi)$.

Je veux d'abord donner une variante affaiblie de la notion de triple de Quillen ( ``Model Category'' de Quillen). \unsure{On|Ce} \unsure{bien|ne} \unsure{sûr|sera} \unsure{pas} une notion tout-à-fait naturelle -- aussi il y aura \uline{deux} variantes affaiblies, duales l'une de l'autre.

Mais avant \unsure{ceci|ça|cela}, un \unsure{sorite|suite} sur les LLP (left lifting property) et RLP (right lifting property). Soit $M$ une catégorie quelconque (le plus souvent, $M$ sera grande). On considère des parties

% slanted
\margin{donner définition en forme, ou renvoi à Quillen}

\[ \Phi,\Psi\subset\mathrm{Fl}(M) \]

et on définit des applications

\[ \boxed{\begin{aligned}\Phi&\longmapsto\Phi_*\\ \Psi&\longmapsto\Psi^*\end{aligned}}\qquad \mathrm{Fl}(M)\longrightarrow\mathrm{Fl}(M) \]

par

\[ \begin{aligned}\Phi_*={}&\{\text{ens. des }p\in\mathrm{Fl}(M)\text{ qui ont \unsure{la|le} RLP p.~r.}\\ &\deleted{\unsure{\text{a}|\text{à}}}\text{à }\Phi\}\\
\Psi^*={}&\{\text{ens des }i\in\mathrm{Fl}\,M\text{ qui ont la LLP p.~r}\\ &\text{à }\Psi\},\end{aligned} \]

et ensuite

\[ \boxed{\begin{aligned}\Phi&\longmapsto\widetilde\Phi\\ \Psi&\longmapsto\overline\Psi\end{aligned}}\qquad \mathrm{Fl}(M)\to\mathrm{Fl}(M) \]

par

\[ \begin{aligned}\widetilde\Phi&=(\Phi_*)^*\\ \overline\Psi&=(\Psi^*)_*\end{aligned} \]

% archive sheet number: 2

\scanpage{4}

\mspage{2}

-- je \unsure{résume} les propriétés les plus immédiates de ces opérations dans une

\uline{Proposition 1} \quad a) $\Phi\to\Phi_*,\ \Psi\mapsto\Psi^*$ sont des applications décroissantes de $\mathrm{Fl}(M)$ dans \unsure{elle-même|elle-m.} pour la relation d'ordre $\subset$, et $\Phi\to\widetilde\Phi,\ \Psi\to\overline\Psi$ des applications croissantes. \deleted{\unsure{On}}

b) \deleted{Appelons $\Phi\unsure{\subset|\in}\Phi$ clos à gauche si \unsure{pour qu'on ait}} On a

\[ \Phi\subset\widetilde\Phi, \]

et pour qu'on ait $\Phi=\widetilde\Phi$, il \unsure{f.} et \unsure{s.} que $\Phi$ soit de la forme $\Psi^*$. On dit alors que $\Phi$ est \uline{clos à gauche}. \deleted{\unsure{\(\Phi\)}} L'ens des parties de $\mathrm{Fl}\,M$ closes à gauche est stable par intersections quelconques$^{*}$, et $\widetilde\Phi$ est la plus petite partie close à gauche de $\mathrm{Fl}\,M$ qui contienne $\Phi$, \inserted{et la plus grande partie $\Phi'$ \deleted{\unsure{de} $\Phi'$ telle que $\Phi'_{*}$} de $\mathrm{Fl}(M)$ telle que $\Phi'_*=\Phi_*$. (ou telle que $\Phi'_*\supset\Phi_*$)}

% slanted
\margin{Donner définition \unsure{bien} en forme.}

% slanted
\margin{$^{*}$ de façon précise $\bigcap_i(\Psi_i^*)=(\bigcup_i\Psi_i)^*$}

Énoncés symétriques pour $\Psi$: \unsure{$\ldots$}

\[ \Psi\subset\overline\Psi, \]

\unsure{On a} $\Psi=\overline\Psi$ \unsure{ss|sss|ssi} $\Psi$ de la forme $\Phi_*$, on dit alors que $\Psi$ est \uline{clos à droite}. L'ens des parties de $\mathrm{Fl}\,M$, closes à droite est stable par intersections quelconques$^*$, et $\overline\Psi$ est la plus petite partie close à droite contenant $\Psi$, et aussi la plus grande partie $\Psi'$ de $\mathrm{Fl}\,M$ telle que $\Psi'{}^*=\Psi^*$. \unsure{(ou} t.q. $\Psi'{}^*\supset\Psi^*$)

% slanted
\margin{$^*$ -- de façon précise, $\bigcap_i(\Phi_{i*})=(\bigcup_i\Phi_i)_*$}

% archive sheet number: 3

\scanpage{5}

\mspage{3}

c) Les applications

\[ \Phi\unsure{\mapstochar\relbar|\mapsto}\Phi_*,\qquad\Psi\mapstochar\relbar\Psi^* \]

établissent des bijections \unsure{inv.|inv} l'une de l'autre entre l'ens \deleted{\unsure{$\Phi$}}des parties $\Phi$ closes à gauche de $\mathrm{Fl}\,M$, et l'ens des parties $\Psi$ closes à droite. \inserted{C'est un anti-isom. d'ens. ordonnés.} On a alors

\[ \begin{aligned}(\Phi\cap\Phi')_{\unsure{*|\alpha}}={}&\text{partie close à droite de }\mathrm{Fl}(M)\\ &\text{``engendrée par }\Phi_*,\Phi'_* \text{'', i.e plus}\\ &\text{petite partie close à droite}\\ &\text{contenant }\Phi_*,\Phi'_*,\ \text{i.e.}\\ &\overline{\Phi_*\cup\Phi'_*}\end{aligned} \]

plus généralement

\[ \left(\bigcap_i\Phi_i\right)_*=\deleted{\text{\unsure{partie}}}\overline{\bigcup_i(\Phi_{i*})} \]

pour tte famille $(\Phi_i)_{i\in\unsure{I|J}}$ de parties closes à gauche de $\mathrm{Fl}\,M$. Formules duales

\[ \left\{\begin{aligned}(\Psi\cap\Psi')^*&=\widetilde{\Psi^*\cup\Psi'{}^*},\\ \left(\bigcap_j\Psi_j\right)^*&=\widetilde{\bigcup_j(\Psi_{j\unsure{*}}^{*})}.\end{aligned}\right. \]

d) Si $\Phi,\Psi\subset\mathrm{Fl}(M)$, alors on a

\[ \Phi\subset\Psi^*\Longleftrightarrow\Psi\subset\Phi_*\qquad\text{on écrit alors }\Phi\longleftrightarrow\Psi, \]

cela signifie que $\Phi$ a la LLP pour $\Psi$, ou que $\Psi$ a \unsure{le|la} RLP pour $\Phi$.

% archive sheet number: 4

\scanpage{6}

\mspage{4}

\uline{Proposition 2}. Soit $\Phi$ une partie close à g. de $\mathrm{Fl}\,M$. Alors $\Phi$ satisfait les conditions a) b) c) ci-dessous :

a) $\Phi$ contient les iso, et est stable par composition, et par cochangement de base (i.e si on a

\diagram{\begin{tikzcd}[ampersand replacement=\&,execute at end picture={\draw[line width=.4pt] (oldx.south west) -- (oldxp.north west) -- (oldy.west) -- (oldxp.east) -- (oldy.north east) -- (oldyp.north east) -- (oldyp.south west);}]
|[alias=oldx]|\deleted{X}\arrow[r] \& |[alias=oldxp]|\deleted{X'}\\
|[alias=oldy]|\deleted{Y}\arrow[u,"\deleted{\unsure{f|j}}"]\arrow[r,"\deleted{i}"'] \& |[alias=oldyp]|\deleted{\unsure{Y'|X'}}\arrow[u,"\deleted{\unsure{f'|j'}}"']
\end{tikzcd}}

\diagram{\begin{tikzcd}[ampersand replacement=\&]
A\arrow[r,"i"]\arrow[d] \& B\arrow[d]\\ A'\arrow[r,"i'",no head]\& B'
\end{tikzcd}}

carré cocartésien \unsure{dans|de} $M$, avec $i\in\Phi$, on a $i'\in\Phi$).

b) $\Phi$ stable par \unsure{facteur direct|facteurs directs} (en tant que partie de l'ens des objets de $\underline{\mathrm{Fl}}(M)$).

c) $\Phi$ \deleted{\unsure{est} $\unsure{\Phi}$} ``stable par limites inductives ordinales'', dans le sens suivant : Soit $\deleted{\unsure{A}}I$ un ens bien ordonné, $(A_i)_{i\in I}$ un \deleted{système} inductif dans $M$, tel que les $u_{\unsure{ji|ij}}:A_i\unsure{-|\to}A_j$ soient dans $\Phi$, et $A_{i_0}\xleftarrow{\sim}\varinjlim_{i<i_0} A_i$ si $i_0\in I$ est un élément limite. Supposons que $A=\varinjlim_I A_i$ existe. Alors $\forall i_0\in I$, la flèche canonique $A_{\unsure{i_0|j_0}}\longrightarrow A$ est dans $\Phi$.

Il y a un énoncé dual pour \unsure{une|les} partie\deleted{\unsure{s}} $\Psi$ close\deleted{\unsure{s}} à droite\unsure{;|'';} \deleted{\unsure{cont.}} elle contient les \unsure{izo|iso}, est stable par composition et par chgt de base, \inserted{par facteurs directs, \unsure{et enfin}} et stable par limites projectives ordinales.

% archive sheet number: 5

\scanpage{7}

\mspage{3}

\deleted{\unsure{$\Phi=\mathcal C$|$\Phi=C$}}

\[ (M,W)\qquad C,\ TC=C\cap W\qquad F=(TC)_*,\ TF=C_* \]

Axiomes de Brown reviennent à

\uline{M0 affaibli} \quad \unsure{sommes amalgamées|somme amalgamée} sur

\diagram{\begin{tikzcd}[ampersand replacement=\&] A\arrow[r,"i"]\arrow[d]\&B\\ A'\& \end{tikzcd}}

\unsure{existe|existent} si $i\in C$ [et j'ajouterais au moins \unsure{;} dualement : \unsure{produit} fibré sur

\diagram{\begin{tikzcd}[ampersand replacement=\&] X\arrow[d,"p"]\&\\ Y\&Y'\arrow[l] \end{tikzcd}}

existe si $p\in F$]

\unsure{M1|M2\,1} n'a pas lieu (\unsure{prolongt relèvement|prolongt, relèvement}) mais

\uline{M2} \uline{affaibl.} \quad a) Tout $f:X\unsure{-|\to}Y$ se factorise \deleted{en} $pi$, avec $i\in C$, et $p\in W$ (on n'exige pas que $p\in\inserted{TF=}C_*$, ni que $C_*\subset W$, ni que $(TC)_*\subset\mathrm{Fib}_W$)

b) $X\longrightarrow e$ se factorise en $X\xrightarrow{i}X'\xrightarrow{p}e$, avec $i\in TC$, $p\in F$.

\uline{M3} Stabilités pour $C$ : contient iso, stable par composition, par cochgt de base

\deleted{\unsure{M4}} On n'a pas besoin de M4, contenu dans M$''$

\uline{M$''$} \quad $C\subset\mathrm{Cof}_W$ [on n'exige pas $F\subset\mathrm{Fib}_W$]

\deleted{\unsure{M'}}, On n'a pas besoin de M$'$, de $TC\subset C\cap W$ (car $TC=C\cap W$), et \deleted{\unsure{on}} n'exige pas $TF=C_*\subset W$.

Ce qui manque, pour avoir une \deleted{\unsure{système}} \inserted{\unsure{str. d}} de Quillen bifaible close à droite, \unsure{exacte} à gauche, est

% slanted
\margin{Je doute que Brown connaisse des cat. à \unsure{cofib.} \unsure{(cf Ch V, 10)|(cf Ch IV, 10)} où \deleted{\unsure{on n'ait|on ait}} où la factorisation \unsure{a)|du n° un} ne puisse se faire \unsure{av|avec} $p\in C_*$ !}

1°) Dans l'axiome de factorisation a), que l'on ait $p\in TF\ (=C_*)$

2°) L'axiome de factorisation $f=pi$, $i\in TC$, $p\in F$, est exigé \unsure{seult} si but/\unsure{src|se|seul} (de $f$) est final,

3°) On n'exige pas $\underset{(TC)_*}{\underset{\Vert}{F}}\subset\mathrm{Fib}_W$, ni $\underset{C_*}{\underset{\Vert}{TF}}\subset W$

% archive sheet number: 6

\scanpage{8}

\mspage{426}

p h y s i q u e : alors même que je me sentais souvent fatigué et sans entrain avant de me mettre au travail - plus ce travail était ``dur'', me faisant manier une lourde pioche ou de grosses pierres disons, plus après coup je me sentais en forme, lourd d'une bonne fatigue. Et il y avait aussi le contact avec des c h o s e s\quad v i v a n t e s : les plantes qu'il fallait soigner ; la terre qu'il fallait préparer pour les accueillir, puis pailler ou biner ; les aliments qu'il fallait préparer et que je mangeais avec autant de plaisir que j'avais eu à préparer le repas ; le chat réclamant sa pitance, et sa part d'affection ; les ustensiles et outils divers aussi, et jusqu'aux cailloux raboteux et souvent mal léchés qu'il fallait tourner et retourner dans tous les sens, afin de les assembler en des murets qui veuillent bien tenir debout...

Effort physique et contact des choses vivantes - ce sont là deux aspects justement qui manquent dans le travail intellectuel, et qui font qu'un tel travail est par nature incomplet, parcellaire, et à la limite, s'il n'est complété et compensé par autre chose, dangereux voire néfaste. C'est la troisième fois, en à peine plus que trois ans, que j'ai eu l'occasion de m'en rendre compte. Il est même devenu bien clair maintenant, que je suis placé devant une échéance draconienne : changer un certain mode de vie, retrouver un équilibre où le pôle yin de mon être, mon corps, ne soit constamment négligé au profit du pôle yang, l'esprit ou (pour mieux dire) la tête - ou sinon, y laisser ma peau dans les toutes prochaines années. C'est ce que mon corps m'a dit, aussi clairement qu'il est possible de le dire ! J'en arrive à présent à un point dans ma vie où le besoin d'une certaine ``sagesse'' élémentaire est devenu une question de s u r v i e, au sens propre et littéral du terme. C'est sûrement une bonne chose- autrement la dite ``sagesse'' se voyait perpétuellement renvoyée aux calendes, au profit de cette espèce de boulimie dans l'activité intellectuelle, qui a été une des forces dominantes dans toute ma vie d'adulte.

Placé devant une échéance aussi claire : ``changer ou crever !'' - je n'ai pas eu à me sonder pour connaître mon choix. C'est ce qui fait que pendant près de quatre mois, j'aie pu, sans jamais avoir l'impression de me faire violence, m'abstenir de toute activité intellectuelle, maths ou pas maths. J'ai su, sans avoir eu à me le dire, qu'à la limite, un jardinier vivant vaut encore mieux qu'un mathématicien mort (ou un ``philosophe'' ou un ``écrivain'' mort, qu'à cela ne tienne !). Avec un peu de malice, on pourrait ajouter : et même mieux qu'un mathématicien vivant ! (Mais ça, c'est encore une autre histoire...).

Je ne crois pas d'ailleurs que je me trouve acculé un jour à une telle

% archive sheet number: 7

\scanpage{9}

\mspage{4}

Les str. de Brown (\unsure{cofibration catégorie|cofibration catégories|cofibration-catégorie} de Brown) permettent de pousser l'alg. homotopique très loin, semble-t-il, en ce qui concerne les structures \deleted{\unsure{impliquées|exactes}} qui tournent autour des suites \unsure{exactes|coexactes}, de suspension, et elles ont \unsure{la|l'air} \unsure{vertu|même} d'être stables par passage de $M$ à $S\backslash M$. Par contre, elles ne sont absolument pas prévues pour la stabilité par passage aux $M/S$, ni pour les suites exactes de cosuspension. \deleted{Mais} Si on \deleted{\unsure{devient}} \unsure{le|la|les} demande, alors l'axiome de factorisation \inserted{b)} $f=pi$ $(\unsure{i|j}\in TC,p\in F)$ est exigé dans \deleted{\unsure{les}} \unsure{sa|se} force (la forme a) restant faible, $f=pi$ avec $i\in C$ et $p\in W$ seulement, pas forcément $p\in F$ et \deleted{\unsure{on}} \unsure{a} fortiori pas \deleted{\unsure{son}} exiger $p\in TF$). Cela ne suffit pas, \deleted{que} faute d'exiger $F\subset\mathrm{Fib}_W$, pour assurer l'existence d'une bonne théorie de fibres $W$-homotopique (alors que l'on en a une, par les seules \unsure{cofibration-catégories|cofibrations-catégories}, pour les \uline{co}fibres $W$-homotopiques).

% archive sheet number: 8

\scanpage{10}

\mspage{427}

situation ``limite'', où j'aurais à renoncer à longue échéance à toute activité intellectuelle, qu'elle soit mathématique ou méditante. Plutôt, la tâche pratique la plus immédiate, la plus urgente dans les années à venir, me semble celle justement d'arriver à un équilibre de vie où les deux types d'activité coexistent au jour le jour, celle du corps et celle de l'esprit, sans que l'une ni l'autre ne devienne dévorante et n'évince l'autre. Je ne me cache pas que c'est bien dans la direction ``esprit'' que se trouvent depuis mon enfance mes investissements les plus puissants, que c'est vers elle aussi que me portent aujourd'hui encore les deux principales passions qui ont continué en ces dernières années à dominer ma vie. De ces deux passions, la passion mathématique et la passion de méditation, il me semble que c'est la première nommée surtout, sinon exclusivement, qui agit comme un facteur de déséquilibre dans ma vie - comme quelque chose qui garderait encore une fâcheuse tendance à ``dévorer'' tout le reste au profit d'elle seule. Ce n'est pas un hasard sûrement, si les trois ``\unsure{épisodes maladie|épisodesmaladie}'' dans ma vie qui ont marqué une situation de déséquilibre, depuis juin 1981, se sont placés en des périodes justement où c'est la passion mathématique qui était sur le devant de la scène.

On pourrait dire que ce n'est pas tout à fait le cas pour ce dernier épisode, survenu en cours de rédaction de Récoltes et Semailles, laquelle constitue une période de réflexion sur moi-même, pour ne pas dire une période de méditation à proprement parler. Mais il est vrai aussi que cette réflexion sur mon passé de mathématicien a été alimentée constamment par ma passion mathématique. Il en a été ainsi surtout dans la deuxième partie, l'Enterrement, il me semble, où la composante égotique de cette passion s'est vue impliquée de façon particulièrement forte et constante. Pourtant, même rétrospectivement, je n'ai pas l'impression qu'à aucun moment, cette réflexion ait pris un rythme, un diapason dévorant, voire démentiel, comme en les deux précédentes occasions où mon corps a été contraint finalement à faire entendre un ``ras le bol !'' sans réplique. Vue séparée du contexte de toute \unsure{'une|une} vie, mon activité intellectuelle depuis un an et demi (depuis la ``reprise'' avec la rédaction de la Poursuite des Champs, suivie par Récoltes et Semailles) apparaît comme se poursuivant à un rythme des plus raisonnables, sans y oublier ni le boire ni le manger (mais parfois quand même, un tantinet, le dormir...). Si elle a fini par déboucher sur un troisième ``épisode santé'' (pour utiliser un euphémisme), c'est sans doute sur le fond de toute une vie marquée par ce sempiternel déséquilibre d'une tête trop forte, imposant son rythme et sa loi à un corps robuste qui a longtemps \deleted{\unsure{encaissé} \unsure{sans} \unsure{broncher} (*).}

% archive sheet number: 9

\scanpage{11}

\mspage{5}

On aura une réciproque partielle dans les catégories accessibles. C'est pour l'énoncé de cette réciproque que j'ai mis à part la condition b) de stabilité par facteurs directs, au lieu de la bloquer avec a).

\uline{Définition 1} \quad a) (Pour mémoire) RLP, LLP, la relation $\Phi\longleftrightarrow\Psi$ entre parties de $\mathrm{Fl}\,M$, (relation non symétrique !).

b) (Pour mémoire) Parties de $\mathrm{Fl}\,M$ closes à gauche, \unsure{close|closes} \unsure{à} droite.

c) Soit $\Phi\longleftrightarrow\Psi$. On dit que le couple $(\Phi,\Psi)$ a la propriété de factorisation de Quillen, si \uline{\unsure{toute|tte}} flèche $f$ de $M$ se factorise en

\[ f=pi,\qquad i\in\Phi,\ p\deleted{\unsure{\Psi}}\in\Psi. \]

On dit que $(\Phi,\Psi)$ est un \uline{couple de Quillen}, si \deleted{\unsure{$\Phi$}} c'est un couple associé par propriétés LLP et RLP $(\Phi\unsure{\leftarrow|\leftarrow\!\shortmid}\Psi)$, \deleted{si elles} et elles sont closes (i.e. $\Phi=\Psi^*,\ \Psi=\Phi_*$), et si elles ont la propriété de factorisation de Quillen\unsure{!|.}

J'énonce dès maintenant le résultat technique essentiel, assurant l'existence de beaucoup de couples de Quillen :

% archive sheet number: 10

\scanpage{12}

\mspage{6}

\uline{Théorème 1} Supposons que $M$ soit \uline{accessible}, et \deleted{que les sommes \unsure{amalgamées|quelconques|inductives}, et par} \inserted{stable par (petites) limites inductives}. (Il doit suffire stabilité par sommes amalgamées, et par $\varinjlim$ filtrantes, voire ordinales\ldots)

\deleted{\unsure{a)}} Soit $\Phi_0$ une \uline{petite} partie de $\mathrm{Fl}\,M$, \deleted{et $\Phi$ une} $\Psi=(\Phi_0)_*$, et $\Phi$ une partie de $\mathrm{Fl}\,M$ telle que

\[ \Phi_0\subset\Phi\subset\Psi^*\ (= \widetilde{\Phi}_0). \]

Alors a) Pour que $\Phi$ soit \deleted{\unsure{stable}} \inserted{close} à gauche, \inserted{i.e. $\Phi=\widetilde{\Phi}_0$}, il f. et s. qu'\unsure{elle|elles} \unsure{satisfasse|satisfassent} les conditions a) b),c), de la prop.~2. Quand il en est ainsi, le couple $(\Phi,\Psi)$ est un couple de Quillen.

b) Supposons seulement que $\Phi$ satisfasse aux conditions a), c) de la prop 2 (à l'exclusion de la condition b) des facteurs directs) -- alors on ne peut affirmer que $\Phi$ soit close à gauche. Néanmoins, $(\Phi,\underset{\Phi_*}{\underset{\Vert}{\Psi}})$ satisfait à la propriété de factorisation de Quillen.

\uline{Remarque} Dans la démonstration, l'existence de \deleted{\unsure{un}} $\Phi_0$\inserted{$\subset\Phi$,} qui joue le rôle d'un petit ens. de ``générateurs\unsure{''} \deleted{pour} (à g.) pour $\Phi$, joue un rôle essentiel. Il y aurait lieu de développer la notion \unsure{des|de} \uline{parties accessibles}

% archive sheet number: 11

\scanpage{13}

\mspage{7}

de $\mathrm{Fl}\,M$, qui impliquerait l'existence d'un \inserted{petit ens.} $\Phi_0\subset\Phi$ tel que $\Phi\unsure{\subset|\in}\widetilde\Phi_0$, et qui aurait la propriété que

\[ \begin{aligned}\Phi\ \text{accessible}&\Rightarrow\Phi_*^{\unsure{*}}\ \text{accessible}\\ \Psi\ \text{accessible}&\Longrightarrow\Psi_{\unsure{,}}^*\ \text{accessible}.\end{aligned} \]

Tout couple \inserted{$(\Phi,\Psi)$} de parties closes \deleted{\unsure{accessibles}} associées \deleted{\unsure{$\Phi=\Psi^*$}} $(\Phi=\Psi^*,\Psi=\Phi_*)$ et accessibles, serait donc un couple de Quillen.

2. \uline{\deleted{\unsure{Triples}} \inserted{Catégories} de Quillen} \uline{faibles}, et \uline{construction d'iceux.}

(Théorème-\unsure{saturation|réduction|situation|solution} conjectural.)

Dans toute la suite, on se donne une catégorie de modèles

\[ \left\{\begin{aligned}&(M,W),\qquad\text{où }W\subset\mathrm{Fl}(M)\\ &W\ \text{faiblement saturé}\unsure{^{*}}:\quad\begin{aligned}&\text{a) } W\text{ contient les iso}\\ &\text{b) Si deux parmi}\\ &f,g,gf\text{ dans }W,\\ &\text{le troisième}\\ &\text{aussi.}\end{aligned}\end{aligned}\right. \]

% slanted
\margin{C'est l'axiome M5 de Quillen}

On veut décrire des données supplémentaires un peu plus faibles que celles d'une cat. de modèles de Quillen. On suppose données des parties

\[ \deleted{\unsure{F|E}}\unsure{C|E},F,TC,TF\quad\subset\quad\mathrm{Fl}(M), \]

avec les propriétés suivantes. Comme axiome préliminaire, je \unsure{pose}

(\unsure{M0|M$_0$}) Stabilité de $M$ par $\varprojlim$ finies et $\varprojlim$ finies.

% archive sheet number: 12

\scanpage{14}

\mspage{8}

M1 \quad Propriété de relèvement-prolongement

\[ \deleted{\unsure{C,F}} \quad C\unsure{\leftarrow|\leftarrow\!\shortmid|\leftrightarrow}TF,\qquad TC\leftarrow F \]

M2 \quad Propriété de factorisation pour $(C,TF)$, $(TC,F)$

[M3] \quad Stabilité de $F$ par composition et par chgt. de base, de $C$ par composition et cochgt. de base, \unsure{.} $F$ et $C$ contiennent les iso.

[M4] \quad Stabilité de $T\deleted{\unsure{C}}F$ \unsure{par} \inserted{composition et} chgt de base, de $TC$ par \inserted{composition et} cochangement de base, et $TF$ et $TC$ contiennent les iso.

\[ \mathrm M'\qquad TC\subset C\cap W,\qquad TF\subset F\cap W \]

\[ \mathrm M''\qquad C\subset\underbrace{\mathrm{Cof}_W}_{W\text{-cofibrations}},\qquad F\subset\underbrace{\mathrm{Fib}_W}_{W\text{-fibrations}} \]

Dans le cas de Quillen (\inserted{Q-}cat. de modèles), M$'$ est renforcé en

\[ TC=C\cap W,\qquad TF=F\cap W \]

et il n'y a pas d'axiome M$''$. L'axiome M$''$ équivaut à la propriété de ``propriety'', \deleted{\unsure{elle, est|et, est}} correspondant aux \inserted{Q-}catégories de modèles \uline{propres}.

\uline{Proposition 3} \deleted{\unsure{M4} \unsure{si l'on} \unsure{une conséquence des autres axiomes} (} \inserted{Supposons les conditions M0 (pour mémoire)} M1, M2, M$'$, M$''$ satisfaites.

Alors en remplaçant $C$, \deleted{\unsure{$F$,}} $TC$, \deleted{\unsure{$F$}} par les \unsure{saturées|saturations} pour les propriétés de stabilité de contenir les iso, \deleted{\unsure{et}} d'être stables par \inserted{composition et par} cochgt de base, \deleted{\unsure{et si on les} \unsure{(les)}} et $F,TF$ par celles de contenir les iso et d'être stables par \inserted{composition et par} chgt. de base, on trouve \deleted{\unsure{$C,TC$,}} $C',TC',F',TF'$ satisfaisant

% archive sheet number: 13

\scanpage{15}

\mspage{9}

à tous les axiomes M1 à M4, M$'$ M$''$.

\uline{Dém.} La prop. 2 nous assure que M1 continue à être satisfaite, \inserted{et} pour M2 c'est trivial, \deleted{\unsure{pour}} M3 et M4 seront valables par construction. D'autre part M$''$ reste valable, car $\mathrm{Cof}_W$ \deleted{\unsure{est}}\unsure{a} les stabilités requises pour $C'$\deleted{\unsure{,$F'$}}, et $\mathrm{Fib}_W$ celles requises pour $F'$. Reste à prouver M$'$, i.e.

\[ TC'\unsure{\subset|\in} C'\cap W,\qquad TF'\subset F'\cap W, \]

et pour ceci il suffit de voir \deleted{\unsure{ceci}} que $C'\cap W$ \unsure{resp.} $F'\cap W$ \deleted{\unsure{ait|ont}}\inserted{\unsure{a}} les stabilités requises pour $TC'$, resp. $TF'$. Cela résulte du

\uline{Lemme}. Soit $C$ une partie de $\mathrm{Cof}_W$. Si $C$ contient les iso, \inserted{il en est de \unsure{$=$} de $C\cap W$. Si de plus $C$} \deleted{\unsure{resp.}} est stable par composition, resp. \deleted{est stable} par cochgt. de base, il en est de \unsure{$=$} $C\cap W$. Énoncé dual pour une partie $F$ de $\mathrm{Fib}_W$.

Cela provient du fait que

\[ \boxed{\mathrm{Cof}_W\cap W=W^{\unsure{\mathrm{univ}|\mathrm{couniv}}}} \]

\inserted{flèches \uline{co-}\uline{universellement} dans $W$}

($\overset{\mathrm{def}}{=}$ ens des flèches $A\unsure{-|\to}B$ telles telles que pour \deleted{\unsure{cochgt de b.|tout chgt de b.}} \inserted{$A\to A'$}, $A'\unsure{-|\to}B'=A'\vee_A B$ est dans $W$

est lui-même stable par composition et par cochgt de base, et une \unsure{$W$-équivalence} \unsure{couniverselle}. Bien sûr, \inserted{si $C\subset\mathrm{Cof}_W$}

\[ C\cap W=C\cap(\mathrm{Cof}_W\cap W)=C\cap W^{\unsure{\mathrm{univ}|\mathrm{un}}}. \]

% archive sheet number: 14

\scanpage{16}

\mspage{10}

Donc il reste à prouver le

\uline{Corollaire \deleted{\unsure{.}} 1} Supposons qu'il existe des parties $C,F,TC,TF\subset\mathrm{Fl}\,M$, satisfaisant \deleted{M1,} M2, M$'{}^{*}$, M$''$. Alors

% slanted
\margin{On n'a pas besoin de \deleted{\unsure{tout}} M$'$, seulement de $TC\subset W$, $TF\subset W$}

\[ \begin{aligned}\mathrm{Cof}_W\cap W&=W^{\unsure{\mathrm{univ}|\mathrm{couniv}}}\\ \deleted{\unsure{\mathrm{Cof}}}\mathrm{Fib}_W\cap W&=W_{\mathrm{univ}}\end{aligned} \]

(flèches \inserted{$f:X\unsure{-|\to}Y$} universellement dans $W$ i.e telles que pour tout $Y'-Y$, $f^{\unsure{\prime\prime|\prime}}:X'=X\times_Y Y'-Y'$ soit dans $W$

Pour voir ceci, montrons p. ex. la deuxième (la première \deleted{\unsure{étant}} se déduira par dualité).

\diagram{\begin{tikzcd}[ampersand replacement=\&] X\arrow[d,"f"']\&X'\arrow[l]\arrow[d,"f'"]\\ Y\&Y'\arrow[l]\end{tikzcd}}

\[ f\in\mathrm{Fib}_W,\ f\in W\ \Longrightarrow f'\in W \]

On a mieux :

\uline{\deleted{\unsure{Corollaire 2}}} \uline{\inserted{Proposition 4}}

\inserted{Supposons que toute flèche de $M$ se factorise en $pi$, avec $i\in W$ et $p\in\mathrm{Fib}_W$. Alors} Si $X,Y$ sont dans $\underline{\mathrm{Fib}}_W/\unsure{S|s}\subset M/S$, alors \deleted{\unsure{toute}} $f:X-Y$ dans $M/S$ qui est dans $W_S$ est dans $W_S^{\unsure{\mathrm{univ}|\mathrm{un}}}$ i.e. elle reste dans $\unsure{W|W_S}$ après \deleted{tout} chgt de base $S'-S$.

\diagram{\begin{tikzcd}[ampersand replacement=\&,column sep=small]X\arrow[rr,"f"]\arrow[dr,"p"',no head]\&\&Y\arrow[dl,"q",no head]\\ \&S\& \end{tikzcd}}

\[ p,q\in\mathrm{Fib}_W \]

Énoncé dual pour les flèches de $S\backslash\underline{\mathrm{Cof}}_W\subset S\backslash M$, \unsure{une|toute} telle flèche qui est dans $W$ est $S$-couniversellement dans $W$, si \deleted{\unsure{tout}} flèche de $M$ se factorise en $qj$, avec $j\in\mathrm{Cof}_W$ et $q\in W$.

% archive sheet number: 15

\scanpage{17}

\mspage{11}

\uline{Dém}

\diagram{\begin{tikzcd}[ampersand replacement=\&,column sep=1.0em,row sep=large]
 X\arrow[drr,"f\in W"]\arrow[ddr]\&{}\&{}\&\overline{X}'\arrow[lll]\arrow[drr,"\overline f'\in W"]\arrow[ddr]\&{}\&{}\&{}\&X'\arrow[llll,"i_X"']\arrow[dr,"f'"]\arrow[ddl]\&{}\\
 {}\&{}\&Y\arrow[dl]\&{}\&{}\&\overline{Y}'\arrow[lll]\arrow[dl]\&{}\&{}\&Y'\arrow[lll,"i_Y"]\arrow[dll]\\
 {}\&S\&{}\&{}\&\overline S'\arrow[lll,"\deleted{\unsure{p|T}}F"']\&{}\&S'\arrow[ll,"TC"', "i"]\&{}\&{}
\end{tikzcd}}

On factorise $S'\unsure{-|\to}S$ en $S'\xrightarrow{i}\overline S'\overset{p}{\unsure{-|\longrightarrow}}S$, avec $i\in\unsure{TC|W}$, $p\in\mathrm{Fib}_W$ par M(2). \deleted{Comme $F\subset\mathrm{Fib}_W$,} \inserted{Donc} $f\in W\Rightarrow\overline f'\in W$. Comme $X,Y\in\mathrm{Ob}\,\underline{\mathrm{Fib}}_W/S$, \inserted{\unsure{ergo}} $\overline X',\overline Y'\in\mathrm{Ob}\,\underline{\mathrm{Fib}}_W/\unsure{\overline S|\overline S'}$, et comme \deleted{\unsure{$p\in W$} \unsure{et} \unsure{$i_X$}} $i:S'-\overline S'$ est dans $W$, donc aussi $i_X,i_Y$, \deleted{\unsure{d'où}} \inserted{on a} $f'\in W$ par M.5 (saturation faible de $W$).

\uline{Commentaire sur Prop 3}

On peut être encore plus exigeant, en demandant dans M3, M4 également la stabilité par facteurs directs, et \deleted{\unsure{lorsque|lors}} Quand seuls M1, M\unsure{2|?}, M$'$, M$''$ sont satisfaits, en saturant $C,TC,F,TF$ également pour cette stabilité supplémentaire. La prop. \unsure{2|?} assure \inserted{encore} que M1 continue à être vérifié, (et \deleted{M2} \inserted{M3, M4 \unsure{trivialemt|trivialement}}), on voit qu'il en est ainsi encore de M$'$, M$''$, à condition qu'on renforce l'axiome de saturation faible \inserted{M5} sur $W$ par

\deleted{Saturation plus forte :} \uline{stabilité de $W$ par facteurs directs} ;

ce qui implique en effet que $\mathrm{Cof}_W,\mathrm{Fib}_W$ \unsure{soient|sont} stables par facteurs directs, ainsi que $W^{\unsure{\mathrm{univ}|\mathrm{couniv}}}$ et $W_{\mathrm{univ}}$.

% archive sheet number: 16

